\subsection{Operational selection of the causal cone}
\label{sec:operational-causal-selection}

The response-selected source Gram space $E$ is three-dimensional. Its unit
directions form a sphere, and adjoining a scalar line gives the quadratic
form $Q(t,x)=t^2-\|x\|^2$ on $\mathbb R\oplus E$. The finite proper carrier
rotations act isometrically on $E$. Under natural feasibility and objective
data with a unique normalized minimizer, the sixty directed seams have
equal weight and select the corresponding translation average on the
conservative record carrier. These are spatial and algebraic results.
The following analytic theorem states an additional condition that selects
an operational causal cone on that same carrier.

A displacement cone $C$ contains zero and is closed under addition and
nonnegative scaling. Suppose it represents a homogeneous operational order
by $p\preceq q$ if and only if $q-p\in C$. Pointedness,
$C\cap(-C)=\{0\}$, expresses antisymmetry. Operational covariance means
that transformations of admissible histories, interventions and readouts
preserve possible influence, with the corresponding linear action on
event displacements. Merely changing a frame or relabeling celestial
directions does not establish this condition. If $C$ is instead the closed
cone generated by finite operations, the conclusion concerns that closure;
it does not assert that every displacement is a finite direct read.

For $u\in E$, define
\[
 X_u(t,x)=(\langle u,x\rangle,tu).
\]
For a unit vector $n$, $B_n(s)=\exp(sX_n)$ is the boost of rapidity $s$:
writing $x=x_\parallel n+x_\perp$,
\[
 t'=\cosh(s)t+\sinh(s)x_\parallel,\qquad
 x'_\parallel=\sinh(s)t+\cosh(s)x_\parallel,\qquad
 x'_\perp=x_\perp.
\]

\begin{theorem}[Finite rotations and one boost direction]
\label{thm:operational-cone-selection}
Let $G\subset\mathrm{SO}(E)$ be finite and irreducible, as for the proper
icosahedral representation. Let $C\ne\{0\}$ be a closed convex pointed
displacement cone invariant under $\operatorname{diag}(1,R)$ for $R\in G$.
Suppose, for one unit $n$, that $B_n(s)C=C$ for every real $s$. Then
\[
 C=\{(t,x):t\geq\|x\|\}
 \quad\hbox{or}\quad
 C=\{(t,x):t\leq-\|x\|\}.
\]
A positive time-axis vector in $C$ selects the first alternative.
It is enough to assume boost invariance for a sequence of nonzero
rapidities tending to zero.
\end{theorem}
\begin{proof}
Irreducibility gives $\operatorname{span}\{Rn:R\in G\}=E$.
Conjugation by $\operatorname{diag}(1,R)$ gives invariance under every
$\exp(sX_{Rn})$. Write $u=\sum_j a_jR_jn$. Since $u\mapsto X_u$ is linear,
\begin{equation}
 \exp(sX_u)=\lim_{m\to\infty}
 \left(\prod_j\exp\!\left(\frac{s a_j}{m}X_{R_jn}\right)\right)^m.
 \label{eq:operational-boost-product}
\end{equation}
Indeed, the one-step product and $\exp(sX_u/m)$ have equal constant and
linear Taylor terms, so their difference in matrix norm is $O(m^{-2})$.
A telescoping sum over $m$ factors bounds the total difference by
$O(m^{-1})$; all intermediate products have a uniform exponential norm
bound. Every finite product preserves $C$, and closedness passes this
inclusion to the limit. Repeating with $-s$ gives equality. Thus every
boost preserves $C$.

The finite average $|G|^{-1}\sum_{R\in G}R$ projects onto the fixed
vectors of $G$, which are zero by irreducibility. Convexity therefore gives
$(t,0)\in C$ whenever $(t,x)\in C$. If $\|x\|>|t|$, boosts parallel to
$x$ can make the time coordinate positive or negative. Averaging these
two images puts both signs of the time axis in $C$, contradicting
pointedness. Hence every nonzero vector in $C$ is causal.

A nonzero future causal vector has positive time coordinate. Its finite
rotation average supplies a positive time-axis vector. Boosts and positive
scaling generate every future timelike vector; closedness adds the null
boundary. A past vector would similarly supply the opposite time axis,
contradicting pointedness. This gives the future cone. Starting with a
past vector gives the past cone. Nontriviality ensures one case occurs.

Finally, if $s_k\ne0$ tends to zero, for every $s\in\mathbb R$ choose
integers $m_k$ with $m_ks_k\to s$. Integer powers and inverses of
$B_n(s_k)$ preserve $C$. Continuity and closedness give invariance under
$B_n(s)$, reducing the last assertion to the preceding proof.
\end{proof}

The theorem uses only two finite representation facts: the boost-axis
orbit spans $E$, and the group average on $E$ vanishes. It does not infer
continuous spatial covariance from the finite carrier action. Its
conclusion is a restriction: rotation covariance alone permits every
cone $\{(t,x):t\geq0,\ \|x\|\leq vt\}$ with $v>0$, and the stationary
time-axis cone. Operational boost covariance excludes those alternatives
except the unit-speed cone in the specified scalar coordinate.

The scalar coordinate requires an additive clock readout
on actual histories. Once that attachment and the normalized boost action
are fixed, the cone slope is fixed relative to them. Constructing the
clock attachment is a separate dynamical requirement. The theorem does not select a layer tick,
read weights, signal amplitudes or a field equation. The count construction
in Section~\ref{sec:source-count-clock} supplies a separate geometric
duration readout on the declared order/count family.

At finite resolution, a bounded event set cannot carry a nontrivial
continuous boost action by permutations. A physical use of the theorem
therefore needs an operational continuum description or controlled
comparisons across a varying family. Finite energy cutoffs and bounded
windows must be transformed or restricted consistently. The existing
algebraic Lorentz and celestial maps alone supply none of these
intervention or clock attachments. The cone classification and product
limit above are analytic proofs, rather than Lean formalizations of
operational covariance.
